% Part: sets-functions-relations
% Chapter: relations
% Section: equivalence-relations

\documentclass[../../../include/open-logic-section]{subfiles}

\begin{document}

\olfileid{sfr}{rel}{eqv}

\olsection{తుల్యతా సంబంధాలు}

ఒక సమితిపై తాదాత్మ్య సంబంధానికి స్వావర్తన, సౌష్ఠవ, సంక్రామక ధర్మాలు
ఉంటాయి. ఈ మూడు ధర్మాలూ ఉన్న $R$ సంబంధాలు చాలా సాధారణంగా కనిపిస్తాయి.

\begin{defn}[తుల్యతా సంబంధం] 
స్వావర్తన, సౌష్ఠవ, సంక్రామక ధర్మాలు ఉన్న $R \subseteq A^2$
అనే సంబంధాన్ని \emph{తుల్యతా సంబంధం} అంటాం.
$A$లోని !!^{element}s $x$, $y$లకు $Rxy$ ఉంటే, వాటిని
\emph{$R$-తుల్యమైనవి} అంటాం.
\end{defn}

తుల్యతా సంబంధాల నుంచి \emph{తుల్యతా వర్గం} అనే భావన వస్తుంది.
ఒక తుల్యతా సంబంధం ప్రవేశాన్ని వేర్వేరు భాగాలుగా ``విడగొడుతుంది''.
ప్రతి భాగంలోని వస్తువులన్నీ ఒకదానితో మరొకటి సంబంధపడి ఉంటాయి;
వేర్వేరు భాగాలకు చెందిన వస్తువుల మధ్య ఆ సంబంధం ఉండదు.
కొన్నిసార్లు ఈ భాగాల గురించి \emph{నేరుగా} మాట్లాడటం ఉపయోగకరం.
అందుకోసం ఒక నిర్వచనాన్ని పరిచయం చేద్దాం:

\begin{defn}\ollabel{def:equivalenceclass}
$R \subseteq A^2$ ఒక తుల్యతా సంబంధం అనుకుందాం.
ప్రతి $x \in A$కూ, $A$లో $x$ యొక్క \emph{తుల్యతా వర్గం} అంటే
$\equivrep{x}{R} = \Setabs{y \in A}{Rxy}$ అనే సమితి.
$R$ కింద $A$ యొక్క \emph{వర్గీకృత సమితి} అంటే
$\equivclass{A}{R} = \Setabs{\equivrep{x}{R}}{x \in A}$,
అంటే ఈ తుల్యతా వర్గాల సమితి.
\end{defn}

తుల్యతా వర్గాలు నిజంగానే $A$ను భాగాలుగా విభజిస్తాయని నిరూపించడం
ద్వారా, ఈ నిర్వచనాన్ని తదుపరి ఫలితం సమర్థిస్తుంది:

\begin{prop}
$R \subseteq A^2$ ఒక తుల్యతా సంబంధం అయితే,
$Rxy$ కావడానికి అవసరమైన మరియు సరిపడే షరతు
$\equivrep{x}{R} = \equivrep{y}{R}$ అవడం.
\end{prop}

\begin{proof}
ఎడమ నుంచి కుడి దిశకు, $Rxy$ అనుకుందాం; $z \in \equivrep{x}{R}$
తీసుకుందాం. నిర్వచనం ప్రకారం $Rxz$. $R$ తుల్యతా సంబంధం కనుక $Ryz$.
(వివరంగా చెప్పాలంటే: $Rxy$ ఉంది, $R$ సౌష్ఠవ సంబంధం కనుక $Ryx$.
అలాగే $Rxz$ ఉంది, $R$ సంక్రామక సంబంధం కనుక $Ryz$.)
అందువల్ల $z \in \equivrep{y}{R}$. సాధారణీకరిస్తే,
$\equivrep{x}{R} \subseteq \equivrep{y}{R}$.
ఇదే విధంగా $\equivrep{y}{R} \subseteq \equivrep{x}{R}$ కూడా.
కనుక మూలకాధారిత సమానత్వం వల్ల
$\equivrep{x}{R} = \equivrep{y}{R}$.

కుడి నుంచి ఎడమ దిశకు, $\equivrep{x}{R} = \equivrep{y}{R}$ అనుకుందాం.
$R$ స్వావర్తన సంబంధం కనుక $Ryy$; అందువల్ల $y \in \equivrep{y}{R}$.
$\equivrep{x}{R} = \equivrep{y}{R}$ అనే ఊహ వల్ల
$y \in \equivrep{x}{R}$ కూడా. కనుక $Rxy$.
\end{proof}

\begin{ex}
తుల్యతా సంబంధాలకు శేషాల ఆధారిత అంకగణితం చక్కని ఉదాహరణను ఇస్తుంది.
ఏ $a$, $b$, $n \in \PosInt$కైనా, $a$ను $n$తో భాగించినప్పుడు
వచ్చే శేషం, $b$ను $n$తో భాగించినప్పుడు వచ్చే శేషంతో సమానమైతే,
అయితేనే $a \equiv_n b$ అంటాం.
(సంకేతాలతో చెప్పాలంటే: ఏదో $k \in \Int$కు $a - b = kn$ కావడం
అనేది $a \equiv_n b$కు అవసరమైన మరియు సరిపడే షరతు.)
ప్రతి $n$కూ $\equiv_n$ ఒక తుల్యతా సంబంధం.
$\equiv_n$ ద్వారా కచ్చితంగా $n$ వేర్వేరు తుల్యతా వర్గాలు ఏర్పడతాయి;
అంటే $\equivclass{\Nat}{\equiv_n}$లో $n$ !!{element}s ఉంటాయి.
అవి: $n$తో శేషం లేకుండా భాగించబడే సంఖ్యల సమితి
$\equivrep{0}{\equiv_n}$; $n$తో భాగించినప్పుడు శేషం $1$ వచ్చే
సంఖ్యల సమితి $\equivrep{1}{\equiv_n}$; \ldots; $n$తో భాగించినప్పుడు
శేషం $n-1$ వచ్చే సంఖ్యల సమితి $\equivrep{n-1}{\equiv_n}$.
\end{ex}

\begin{prob}
ప్రతి $n \in \PosInt$కూ $\equiv_n$ తుల్యతా సంబంధమనీ,
$\equivclass{\Nat}{\equiv_n}$లో కచ్చితంగా $n$ సభ్యాలు ఉంటాయనీ చూపండి.
\end{prob}

\end{document}
